\section{An alternative rational verification procedure}
\label{app:floor-certificate}

This appendix proves a conditional alternative to the exact-array interval
verification used in Proposition~\ref{prop:finite-certificate}.  If the
rational procedure below passes the stronger tables of
Section~\ref{sec:finite}, its roundoff allowances leave the exact quadrature
arrays within the proposition's looser bounds.  The complete floor-rounded
matrix procedure is not executed by the supplied programs; the selected
finite certificate uses direct Arb enclosures of the exact arrays.

We use $\epsilon=10^{-80}$ and the decimal brackets $p_*,b_*$ from
\eqref{eq:rational-pi-b} in Appendix~\ref{sec:arithmetic}.  We first define
the rounding and prove its error bounds.  The complete listings then specify
each finite loop, keeping all arithmetic not explicitly rounded exact.

\subsection{Componentwise rounding and its error}

For real $x$, let $R(x)=\epsilon\lfloor x/\epsilon\rfloor$.  Apply $R$
separately to the real and imaginary components of a complex number and
entrywise to a matrix.  For complex rational $w$, define
\begin{equation}\label{eq:rational-exponential}
\begin{gathered}
 u=R(w/4096),\qquad d_0=1,\qquad
 d_j=R(d_{j-1}u/j)\quad(1\le j\le60),\\
 x_0=\sum_{j=0}^{60}d_j,\qquad
 x_{j+1}=R(x_j^2)\quad(0\le j<12),\qquad e_*(w)=x_{12}.
\end{gathered}
\end{equation}
This uses a degree-$60$ polynomial with $61$ terms followed by twelve
squarings.  Every quantity is rational.  In the floor prescription,
$\pi,b,\sqrt3$ are replaced by $p_*,b_*,2b_*$, cosine is evaluated as
$\Rea e_*(iw)$ at the corresponding rational phase $w$, and $B,h$ remain
exact.

In addition to the floors internal to $e_*$, the prescription rounds the
quadrature points and weights, mass entries, and values of $z$ and $\lambda$.
It rounds the transformed zeroth moment factor and every higher moment
factor after its divided-power update.  It also rounds the projected jet
entries, the entries of $Y$, every product in the power and partial-sum
updates, and both final positive-node coefficient matrices.  The extra
coordinates and all other operations are exact rational arithmetic; in
particular, the products used to form the Bernstein and tail rows incur no
further floors.  The listings below fix the order of these operations.

A superscript $\#$ denotes the output of this floor prescription.  Its
error is measured against the corresponding exact quadrature array of
Section~\ref{sec:finite}: for example,
$S^\#_{\mathcal F\mathcal F}$ is compared with
$\widetilde S_{\mathcal F\mathcal F}$, $S^\#_{\mathcal F e}$ with
$\widetilde S_{\mathcal F e}$, and $Y^\#,X_i^\#$ with $Y,X_i$.
The error from quadrature to the exact spectral integrals is a different
quantity, already bounded in Lemma~\ref{lem:quadrature-error}.

\Needspace{6\baselineskip}
\begin{lemma}[Roundoff bounds for rational verification]
\label{lem:rational-roundoff}
The floor prescription below has the following properties.
\begin{enumerate}
\item If $|w|<1700$ and $\Rea w<7$, then $|e_*(w)-e^w|<10^{-68}$.
Every exponential argument in the prescription lies in this domain.
\item A computed single-column quadrature moment of order at most $42$,
scaled by $y^a/a!$ with $|y|\le3/2$, differs from the exact quadrature
moment by less than $10^{-48}$.  The initial operator-norm errors in
$S^\#_{\mathcal F\mathcal F}$, $S^\#_{\mathcal F e}$, and $Y^\#$ are
less than $10^{-44}$.
\item Suppose that either the exact or the rounded computation satisfies
the following envelopes: every pre-update power in the nine steps has norm
below $3$, each intermediate partial-sum matrix has norm below $34$, and
every combined pair of corresponding $X_1,X_2$ columns has norm below $5.6$.
Then the combined error in a pair of corresponding columns of
$X_1^\#,X_2^\#$ is less than $10^{-29}$.
Every entry of the Bernstein rows $L_{i,d}^\#$ differs from the
corresponding entry of $L_{i,d}$ in \eqref{eq:bernstein-rows} by less than
$10^{-16}$.  Every lower evaluation in \eqref{eq:bernstein-certificate}
and every tail expression in \eqref{eq:tail-certificate} has error less
than $10^{-6}$.
\end{enumerate}
The stagewise envelopes in the last item follow from the stronger power and
column tables, whether those tests are established on the exact arrays by
intervals or on the rounded arrays by the prescription below.
\end{lemma}

\begin{proof}
We control elementary functions first, then a single moment, then the
matrix products, and finally the scalar lower evaluations.  This order
keeps the error due to the coefficient lists separate from the error in
the primitive moments.

\paragraph{Exponential approximation.}
A componentwise floor changes a complex number by less than $2\epsilon$.
For $r=w/4096$, both $|r|$ and $|u|$ are below $.416$.  If $\delta_j$
bounds $|d_j-r^j/j!|$, then
\[
 \delta_0=0,\qquad
 \delta_j\le\frac{.416}{j}\delta_{j-1}
 +\frac{2\epsilon\,.416^{j-1}}{j!}+2\epsilon.
\]
The first term propagates the previous error, the second changes $r$ to
$u$, and the last is the new floor.  Summing through $j=60$ and adding
the exponential tail gives
\[
 |x_0-e^{w/4096}|
 \le\sum_{j=0}^{60}\delta_j+\frac{e^{.416}.416^{61}}{61!}
 <1.263\cdot10^{-78}<4\cdot10^{-78}.
\]
The factorial tail alone is less than $4\cdot10^{-107}$.  These are finite
rational comparisons after using $e^{.416}<e<2719/1000$.

At squaring step $j$, the exact value has modulus at most
$e^{7/2^{12-j}}$.  While the accumulated error is below one, its
amplification factor is at most $2e^{7/2^{12-j}}+1$.
Since $2e^a+1\le3e^a$ for $a\ge0$, the product of all twelve factors
is at most $3^{12}e^7$.  Including all twelve new floors gives
\begin{equation}\label{eq:exponential-roundoff-budget}
 (4\cdot10^{-78}+24\epsilon)3^{12}e^7
 <2.472\cdot10^{-69}<10^{-68}.
\end{equation}
The same bound is below one at every earlier stage, closing the hypothesis
used in the amplification estimate.

The arguments used by the prescription have ample room in the stated domain.
Density phases have modulus at most $\pi\cdot168$ before the rational
perturbations.  The weight cosines have phase modulus less than
$2\pi\cdot191$.  Transformed phases have modulus at most
$\pi(44/15)168<1549$ and nonpositive real part by
Lemma~\ref{lem:kernel-bounds}.  They remain in the left half-plane after
substitution: on the real integration interval, $\Ima z>.187/\pi$,
whereas the perturbation of $z$ below is less than $10^{-65}$.
The only positive real scalar phases are $K\le6$ in the support at $m=0$.
Every other finite scalar phase is nonpositive with modulus at most $48$.
Infinite endpoints are evaluated by their limits, not passed to $e_*$.
Thus every call satisfies $|w|<1700$ and $\Rea w<7$.

\paragraph{Primitive masses and moments.}
The intermediate absolute error allowances are
\begin{equation}\label{eq:primitive-roundoff-budgets}
\begin{array}{c|c}
\text{quantity}&\text{absolute error bound}\\ \hline
 \tau_{jv}&10^{-67}\\
 o_v&10^{-65}/384\\
 z(t),\ \lambda(t)&10^{-65}\\
 \text{one density value}&10^{-60}\\
 \text{one mass column, summed absolute error}&10^{-59}\\
 F_0(t),\ G_0(t)&10^{-60}\\
 F_a(t),\ G_a(t),\quad0\le a\le42&10^{-52}\\
 Q_n,\ D_n&10^{-75}.
\end{array}
\end{equation}
The comparison values here use the exact quadrature points and the exact
constants $\pi,b$.  We justify each stage rather than infer it from decimal
agreement.

For a weight phase, replacing $\pi$ changes the phase by less than
$382\epsilon$.  Add the $10^{-68}$ exponential error, and use
$2\sum_{a=1}^{191}(4a^2-1)^{-1}<1$.  Including the final floor gives
the stated bounds for $\tau$ and $o$.  On the real line, $|t+ih|\ge h$;
the derivative bounds
\[
 |z'(t)|\le B/h^2<9,\qquad
 |\partial_t\lambda(t)|<8,\qquad |\partial_b\lambda(t)|<4
\]
then give the errors for $z$ and $\lambda$.  The density phase errors
are below $10^{-63}$ because $n\le168$.  Inserting them in the six-term
density formulas, using $\sum_{l>0}|P_l|=10$, gives the density allowance
$10^{-60}$. The weighted density errors sum to at most $10^{-60}$ because
the six positive quadrature rules have total weight one. Weight errors,
density errors, the $2311$ componentwise mass floors, and atom coefficient
errors together contribute less than $10^{-59}$
to the sum of absolute errors in one mass column.

Each exact mass column has sum of absolute values less than $230$.  Indeed,
the pointwise real-interval density bounds are
$2\pi^2\sum_{l>0}t_l|P_l|<106$ for a $c_n$ column and $20\pi<63$ for a
$d_n$ column; the constant column has atomic variation $25$.  Positivity
and total density weight one from \eqref{eq:positive-quadrature-weights}
make $230$ a common bound.

The direct zeroth kernel phase error is below $10^{-63}$.  The transformed
zeroth phase error is below $6\cdot10^{-63}$, and multiplication by
$|\lambda|<3$ still leaves error below $10^{-60}$.  The divided-power
arguments $i\pi yt$ and $i\pi yz$ have modulus below $15$ and error below
$10^{-63}$.  To see explicitly how these errors enter a higher moment, write
either exact factor as $H_a=H_0v^a/a!$ and its rounded update as
$H_a^\#=R(H_{a-1}^\#v^\#/a)$.  Here
\[
 |H_0|<3,\qquad |H_0^\#-H_0|<10^{-60},\qquad
 |v|,|v^\#|<15,\qquad |v^\#-v|<10^{-63}.
\]
Changing the initial value costs at most $e^{15}10^{-60}$, and changing
$v$ costs at most $3e^{15}10^{-63}$, by the derivative of $v^a/a!$.
A floor introduced at step $j$ is multiplied by
$(v^\#)^{a-j}j!/a!$.  Since $j!/a!\le1/(a-j)!$, all the floors together
cost at most
\[
 2\epsilon\sum_{j=1}^a15^{a-j}\frac{j!}{a!}\le2\epsilon e^{15}.
\]
Thus for $0\le a\le42$ the total factor error is bounded by
\[
 |H_a^\#-H_a|
 <e^{15}\bigl(10^{-60}+3\cdot10^{-63}+2\epsilon\bigr)<10^{-52}.
\]
The last comparison follows already from $e<3$.
Expanding the moment difference against the rounded mass column, whose norm
is below $230+10^{-59}$,
also retains the product of the two primitive errors.  The total is at most
\[
 (230+10^{-59})10^{-52}+3e^{15}\cdot10^{-59}<10^{-48}.
\]
Finally, \eqref{eq:node-constants} and \eqref{eq:rational-pi-b-error}
give errors below $10^{-75}$ in $Q_n,D_n$.  Jet projection uses
$Q_n>1.75$, $|D_n|<2.12$, and $168$ output coordinates.  Including its
floors, the resulting initial operator errors are less than $10^{-44}$;
the same estimate applies to $Y^\#$.

\paragraph{Powers and coefficient columns.}
A floor on a real matrix with $168$ rows contributes at most
$168\epsilon<10^{-77}$ in operator norm.  Let $e_r$ bound the error in
$U$ before step $r$, and let $v_r$ bound the error in either partial-sum
matrix after $r$ steps, with $v_0=0$.  The assumed stagewise envelopes give
$\norm{U}<3$ and partial-sum norms below $34$.  Comparing products and
then adding the floor gives
\begin{equation}\label{eq:matrix-roundoff-recurrence}
 e_0<10^{-44},\qquad
 e_{r+1}\le7e_r+10^{-77},\qquad
 v_{r+1}\le5v_r+34e_r+10^{-77}.
\end{equation}
More precisely, the square error is at most $(6+e_r)e_r+10^{-77}$,
and the partial-sum coefficient of $v_r$ is at most $4+e_r$.
Induction gives $e_r<1$ and validates the displayed recurrence.  Its
finite rational solution yields
\[
 \max_{0\le r\le8}e_r<5.765\cdot10^{-38}<10^{-36},
 \qquad v_9<6.529\cdot10^{-36}<10^{-34}.
\]
If the coarse norms are known for the rounded rather than exact factors,
exchange their roles in the product difference.  The same two inequalities
follow.  This is why the argument is not circular and does not require a
rounded matrix inverse.

The largest combined column norm of $E_1,E_2$ is $3.3505$.
The formulas for $Y$ give $\norm{Y}<2$, and the transform bound gives
the sharper data estimate
$2+(4000+2\cdot10^{-27})3.3505<13405$.
Its rational perturbation is less than $10^{-43}$. Thus, for every column $j$,
$\norm{(g_1)_j}_1+\norm{(g_2)_j}_1<20000$ for both the exact and rounded data.
Applying the two sum/difference formulas in
\eqref{eq:finite-columns}, including the $168$ real floors in each of
the two final columns, gives the complete bound
\[
 40000v_9+72\cdot10^{-43}+336\epsilon
 <2.612\cdot10^{-31}<10^{-29}.
\]
The appended extras are exact.  In particular, the term $336\epsilon$ is
part of the column error even though it is negligible numerically.

\paragraph{Bernstein and lower-evaluation errors.}
A scaled moment against a unit list column has absolute value at most
$230\cdot3e^{15}$.  Also $|y|^{-\nu}\le4$ and
$0\le\omega_{dr}\le1$.  The contribution from the coefficient-list error
to one Bernstein row entry is therefore at most
\[
 41\cdot4\cdot230\cdot3e^{15}\cdot10^{-29}
 <3.700\cdot10^{-18}.
\]
For this decomposition the list difference is evaluated with exact moments;
the primitive difference is evaluated against the rounded columns.
The primitive moments therefore contribute separately.  The relevant combined
column norm is below $6$, either directly from the rounded cutoffs or
from the exact cutoff $5.6$ and the column error.  Thus their contribution
is at most
\[
 41\cdot4\cdot6\cdot10^{-48}<10^{-44}.
\]
The sum of these two contributions is below $10^{-16}$, as claimed.
This explicitly includes the primitive-moment addition.

For completeness, the subsequent minimizations do not amplify that bound
to anything close to $10^{-6}$.  The same moment bound with combined
column norm $6$ shows that each exact Bernstein row entry is below $10^{14}$
in modulus.  If its entry error is $\delta=10^{-16}$, then for
$x\in\{1/K,1/J\}\subset[0,1]$ the error in $a_j+xb_j$ is below $2\delta$.
A finite endpoint exponential has error below $10^{-68}$ and modulus at
most one; the infinite-endpoint values are exact.  Each product in
\eqref{eq:box-lower} consequently has error below
$2\delta+(2\cdot10^{14}+2\delta)10^{-68}$.  There are four such products
and one unmultiplied term.  A minimum is nonexpansive in the maximum norm,
so the total row lower-evaluation error is less than $2\cdot10^{-15}$.

Only the target supports at $m=0$ and $m=3$ use a rounded exponential.
At $m=0$, $|Ky|\le3$, $1/K<1$, and
$\sum_{r=0}^{40}3^r/r!<e^3<27$, so their total weighted error is below
$27\cdot10^{-68}$.  At $m=3$, $J/2\le3$ gives error below
$3\cdot10^{-68}$.  The $m=1$ support is rational and exact.  These
additions still leave a bound far below $10^{-6}$.

For a tail row from $X_i$, the entry error is below $10^{-29}$ and its
entry modulus is below $6$.  The same is true for the row obtained by
summing the low $d_n$ rows, because the combined column norm controls that
sum.  The preceding five-term argument gives lower-evaluation error below
$2\cdot10^{-28}$ for each such row.  Taking $\min(0,\cdot)$ does not
increase it.  In \eqref{eq:tail-certificate} there are one constant row,
one summed row, and two rows for each of $46$ low nodes, hence $94$ terms.
Their nonnegative reciprocal factors are at most one since $s_0-n\ge1.5$.
The tail error is therefore below $94(2\cdot10^{-28})<2\cdot10^{-26}$,
in particular below $10^{-6}$.
\end{proof}

\subsection{The complete floor-rounded finite prescription}

We now specify the alternative rational computation.  Every loop has a
finite inclusive range.  The five shell and two list indices are one-based;
quadrature and derivative indices start at zero as indicated.  Products
are ordinary matrix or row products, without conjugation.  An operation
not marked $R$ is exact rational arithmetic.  The imaginary unit is
\texttt{i}; a list subscript never denotes that unit.  The symbols
$P_l,Q_n,D_n$ mean the formulas in \eqref{eq:periodic-coefficients} and
\eqref{eq:node-constants} after the rational substitutions above.

The input order is all $c_n$ for increasing $n\in\mathcal F$, then all
$d_n$, then $c_0,d_0,C$.  The $2311$ mass rows are the $6\cdot384$
quadrature points in lexicographic order followed by the seven atoms.

\Needspace{12\baselineskip}
\begin{lstlisting}[basicstyle=\ttfamily\footnotesize,columns=fullflexible,
breaklines=true,frame=single,caption={Floor-rounded verification: masses and moment factors},label=alg:finite-verification]
eps = 1/10^80
R(x) = eps * floor(x/eps)  # separately on complex components
# pstar,bStar are the two printed 80-digit decimal rationals.
# e_star uses the degree-60 Taylor polynomial and 12 squarings above.

F_nodes = [n for n=1..168 if n mod 12 in {0,1,3,4,7,9}]
t_l = l/6 for l=0..6
for v=0..383:
    a0 = (v+1/2)/384
    o[v] = R((1-2*sum(Re(e_star(2*d*i*pstar*a0))/(4*d*d-1)
                          for d=1..191))/384)
    for j=1..6:
        tau[j,v] = R((2*j-1+Re(e_star(i*pstar*a0)))/12)

W = zero complex matrix with 2311 rows and 171 columns
for j=1..6:
  for v=0..383:
    t = tau[j,v]
    for n in F_nodes followed by [0]:
      for l=j..6:
        A[l] = (o[v]/6)*2*pstar*e_star(-i*pstar*t*n)
                            *P[l]*e_star(i*pstar*t_l*n)
      W[(j,v),c(n)] = R(-pstar*sum((t_l-t)*A[l] for l=j..6))
      W[(j,v),d(n)] = R(i*sum(A[l] for l=j..6))
for j=0..6:
    W[atom(j),C] = R((1 if j==0 else 2)*P[j])

function factors(m,y,end_degree):
    points = all tau[j,v], followed by t_0,...,t_6
    for each point t:
        z[t] = R(-B/(t+i*h)-i*h)
        lam[t] = R(i/(bStar*(t+i*h)))
        F[0,t] = e_star(i*pstar*t*m)
        G[0,t] = R(lam[t]*e_star(i*pstar*z[t]*m))
        for a=1..end_degree:
            F[a,t] = R(F[a-1,t]*i*pstar*t*y/a)
            G[a,t] = R(G[a-1,t]*i*pstar*z[t]*y/a)
    return F,G
\end{lstlisting}

These moments are scaled by $y^a/a!$, which explains the division by
$y^\nu$ in the later half-gap rows.  The floor convention can itself be
implemented with integers.  Put $S=10^{80}$.  For $w=w_r+iw_i$, initialize
$u_r=\lfloor Sw_r/4096\rfloor$ and
$u_i=\lfloor Sw_i/4096\rfloor$.  Represent a rounded complex number by
the two integers obtained after multiplication by $S$.  The Taylor update is
\[
 (d_r,d_i)\longmapsto
 \left(
 \left\lfloor\frac{d_ru_r-d_iu_i}{Sj}\right\rfloor,
 \left\lfloor\frac{d_ru_i+d_iu_r}{Sj}\right\rfloor
 \right),
\]
with both right sides evaluated before replacement.  The squaring update is
\[
 (x_r,x_i)\longmapsto
 \left(
 \left\lfloor\frac{x_r^2-x_i^2}{S}\right\rfloor,
 \left\lfloor\frac{2x_rx_i}{S}\right\rfloor
 \right).
\]
Every floor is toward negative infinity, including for negative numerators.
These formulas are exactly \eqref{eq:rational-exponential}.

\Needspace{12\baselineskip}
\begin{lstlisting}[basicstyle=\ttfamily\footnotesize,columns=fullflexible,
breaklines=true,frame=single,caption={Floor-rounded verification: projected matrices and columns},label=alg:finite-columns]
Ssharp = zero matrix with 168 rows and 171 columns
Ysharp = zero matrix with 168 rows and 10 columns
for m in F_nodes:
    F,G = factors(m,1,1)
    a = Re(G[0,:] * W)
    d = Re(G[1,:] * W)
    Ssharp[c(m),:] = R(-a/Q_m)
    Ssharp[d(m),:] = R((D_m*a-d)/Q_m)
for j=1..5:
    n = [1,3,4,7,9][j]
    Ysharp[d(n),j]   = R(-1/Q_n)
    Ysharp[c(n),j+5] = R(1/Q_n)
    Ysharp[d(n),j+5] = R(-D_n/Q_n)

E = zero matrix with 6 rows and 10 columns
E[:,1] = [.567,.250,-.0074,0,-.209,.0096]
E[:,6] = [-.976,.408,.0115,1,.938,-.017]
E1 = first three rows of E
E2 = last three rows of E
U = positive-coordinate columns of Ssharp
vp = identity matrix of size 168
vm = identity matrix of size 168
for r=0..8:
    vp = vp + R(vp*U)
    vm = vm + R((-1 if r==0 else 1)*vm*U)
    record the maximum absolute column sums of U,vp,vm
    if r<8: U = R(U*U)
# Compare every recorded value with the power table, row r.

Se = last three columns of Ssharp
g1 = Ysharp + Se*E2
g2 = Se*E1
X[1] = stack(R((vp*(g1+g2)+vm*(g1-g2))/2), E1)
X[2] = stack(R((vp*(g1+g2)-vm*(g1-g2))/2), E2)
# Test every individual and same-column combined absolute column sum
# against the corresponding bound in the column table.
\end{lstlisting}

The product identities in the proof of
Proposition~\ref{prop:finite-certificate} give the prescribed $512$-term
series.  At each stage, the entries of \texttt{vp} and \texttt{vm} are
integer multiples of $\epsilon$.  Hence
$V+R(VU)=R(V(I+U))$, with the analogous identity for the initial minus
sign.  The listing and the rounded product update used in
\eqref{eq:matrix-roundoff-recurrence} are therefore the same operation.

\Needspace{12\baselineskip}
\begin{lstlisting}[basicstyle=\ttfamily\footnotesize,columns=fullflexible,
breaklines=true,frame=single,caption={Floor-rounded verification: every half-gap and tail inequality},label=alg:finite-signs]
nodes5 = [1,3,4,7,9]
endpoints = [2.36,2.65,3.2,4,6,infinity]
sigma = [-1,1]

function low(row,K,J):
    exK = [e_star(-K*(n-1)) for n in nodes5]
    exJ = ([e_star(-J*(n-1)) for n in nodes5]
                   if J<infinity else [1,0,0,0,0])
    answers = []
    for x in [1/K, (1/J if J<infinity else 0)]:
        v = row[1..5] + x*row[6..10]
        answers.append(v[1] + sum(min(v[a]*exK[a],v[a]*exJ[a])
                                                     for a=2..5))
    return min(answers)

low_nodes = [0] followed by all n in F_nodes with n<=88
for m in low_nodes:
    # Successor and predecessor are in {0} union the positive nodes.
    ys = [(successor(m)-m)/2]
    if m>0: append (predecessor(m)-m)/2 to ys
    for y in ys:
        nu = (0 if m==0 else 2)
        F,G = factors(m,y,40+nu)
        for fn=1..2:
            for a=0..40:
                l[a] = sigma[fn]/y^nu * Re(
                    F[a+nu,:]*(W*X[fn])
                    +G[a+nu,:]*(W*X[3-fn]))
            for consecutive K,J in endpoints:
                tc = zero row of length 41
                if m==0:
                    tc[r] = e_star(K)/K * (-K*y)^r/r! for r=0..40
                if m==1:
                    tc[r] = K*(-K*y)^r/(r+2)! for r=0..40
                if m==3 and J<infinity:
                    tc[0] = (J/2)*e_star(-13*J/6)
                for d=0..40:
                    lsum = zero row of length 10
                    tsum = 0
                    w = 1
                    for r=0..d:
                        lsum = lsum + w*l[r]
                        if fn==1: tsum = tsum + w*tc[r]
                        if r<d: w = w*(d-r)/(40-r)
                    value = low(lsum,K,J) + tsum
                    # Strictly compare with the m-group and K row
                    # of the half-gap table, for every value.

s0 = 89.5
for consecutive K,J in endpoints:
    for fn=1..2:
        signedX = sigma[fn]*X[fn]
        neg(row) = min(0,low(row,K,J))
        value = low(signedX[C,:],K,J)
        value += neg(sum(signedX[d(n),:] for n in low_nodes))/s0
        for n in low_nodes:
            value += neg(signedX[c(n),:])/(s0-n)^2
            value += n*neg(signedX[d(n),:])/(s0*(s0-n))
        # Strictly compare with the tail column for this K row.
\end{lstlisting}

The three references to tables in these listings mean, respectively,
Tables~\ref{table:power-bounds}, \ref{table:column-bounds}, and
\ref{table:sign-bounds}.  Every comparison is strict and
rational.  There are $37310$ half-gap comparisons and ten tail comparisons;
a reported extremum is not a substitute for testing each index.

If the floor computation passes all these tests, the roundoff lemma transfers
them to the looser exact-array bounds in
Proposition~\ref{prop:finite-certificate}.  For example,
$.00116-10^{-6}>.001$ and $.0022-10^{-6}>.0021$.
The column cutoffs $2.20,5.40$ remain below $2.3,5.6$ after the column
error, and the power errors are much smaller than the gaps from $.0031$
to $.004$ and from $32.12+1.85$ to $36$.

